\section*{The Propers: Detailed Commentary}

\subsection*{The petitioning assembly\enspace\elementref{Entrance, Collect}}

Psalm 27 lets need and desire occupy the same prayer: the threatened speaker asks for God's hearing and help while seeking one thing, life in God's presence. At the Entrance only the plea is appointed; the house-of-God verse belongs to Communion option A. Their common psalm creates a genuine textual arc across the Entrance and Communion A.

Augustine's \emph{Exposition on Psalm 27} receives the prayer as the voice of Christ's body learning not to confuse delayed relief with divine abandonment. His treatment keeps petition and worship together: the suppliant asks to be heard, but the desired good is finally the Lord himself rather than merely escape from enemies. His reception of the psalm is direct and bears on the arc the two antiphons trace. Entrance and Communion A are not two unrelated snippets: the cry for help and the desire to dwell belong to one scriptural act of confidence. Because Communion A remains an unselected branch, that fuller arc is available rather than enacted by default.

The Collect's Latin identity places God himself as the strength of those who hope in him. Its anthropology is neither pessimistic nor voluntarist. Human willing and acting are real, but frail; grace enables the creature to intend and perform what pleases God.

\subsection*{A treasured people under the Lord of the harvest\enspace\elementref{Exodus, Psalm, Matthew}}

Exodus 19 follows liberation and precedes ratification. Israel is not selected because it achieved holiness; God first carried and gathered the people. “Possession,” “kingdom of priests,” and “holy nation” therefore name gift and task together. Catholic reception applies this vocabulary to the Church in Christ, while Exodus continues to name Israel's election within its own covenant history.

The syntax of the oracle is conditional without being transactional. God's prior action---bringing Israel out and bearing the people---precedes the summons to hear and keep covenant. The people's priestly and holy identity is consequently received within a relationship God initiated, not purchased as a reward; it is a vocation for the sake of a world already named as his possession. The appointed boundary ends before the people's response and the covenant rite: 19:2--6a announces the vocation but does not contain the whole Sinai event.

Augustine's \emph{Letter 55}, section 30, does not expound every phrase of Exodus 19:2--6a. It instead receives Sinai in a chronology from Passover to the giving of the Law and sets that fifty-day pattern beside Christian Pentecost and the gift of the Spirit. Augustine's chronology illuminates without expounding: it places the covenant scene within a Paschal--Pentecostal economy, while the narrated election is Exodus's own. The major accessible English Greek- and Latin-Father collections and Latin lead corpora yield no comparably sustained early exposition of the exact appointed verses, and none of the Fathers read in them supplies an exposition behind the later ecclesial language.

The \emph{Apostolic Constitutions} II.25 also reuses the priestly-people vocabulary in describing an ordered Christian assembly. The work attests early ecclesial appropriation but does not comment on Exodus directly. Together with Augustine's chronology, it records reception; neither treats the exact appointed unit at length.

Psalm 100 turns the covenant declaration into congregational response. The sheep do not make their shepherd, and the people do not create their own belonging. Yet service is active and joyful. The psalm holds together creaturehood, covenant identity, worship, and enduring divine fidelity.

Augustine's direct exposition follows the psalm's widening address to all lands and refuses a merely local possession of praise. He lingers over joyful service, God's making of the people, and the flock's dependence on its shepherd. He reads the psalm within the Church of his own day, not within its original Israelite occasion, and that reading helps explain why the response does more than repeat Exodus. The first reading announces a people made God's treasured possession; the psalm teaches that people how to answer---with glad worship, acknowledged creaturehood, and trust in mercy extending beyond a single generation.

Matthew's Jesus first sees. Compassion is not a decorative feeling but the source of prayer, naming, authority, and sending. The mixed images of shepherdless flock and abundant harvest both keep agency with Christ. Chrysostom notes that prophetic labor has preceded the apostles; they reap what grace has prepared, while Christ is expressly the harvest's Lord. Mission therefore resists both despair, as though the field were empty, and triumphalism, as though the workers owned it.

Chrysostom's \emph{Homily 32 on Matthew} follows the passage's sequence closely. The crowd's condition moves Christ to compassion; the harvest saying makes him not one laborer among others but the Lord who sends; and the named Twelve receive power before they exercise it. Chrysostom also notices that the harvest is already prepared by labor before the apostolic mission. Aquinas's later commentary on John 4:37--38 offers an illuminating parallel rather than direct Matthew commentary: prophets and Christ prepare what apostles reap, so missionary fruit cannot be credited to the visible reaper alone. Read together, the witnesses prevent an institutional history in which apostolic action begins from nothing.

The passage also orders prayer and office. Before the Twelve are named, disciples are told to ask the harvest's Lord to send laborers; the answer to that prayer takes the form of persons called by name and given authority. Prayer is therefore not the pious alternative to institutional mission, and office is not a self-originating mechanism. Each remains answerable to Christ's perception of the harassed crowd. The list itself includes differences and future failures that the selected verses do not narrate: Matthew the tax collector stands with Simon the Zealot, and Judas is named as the one who betrayed Jesus. Authority is genuinely conferred without canonizing every later act of its bearer.

The verbs of 10:7--8 keep word and deed inseparable. Proclamation of the Kingdom is accompanied by healing, cleansing, and liberation; none becomes a commodity because all authority is received. The verbs promise no guaranteed technique for miraculous cure, yet they leave no room to spiritualize compassion into speech that ignores bodies. Chrysostom's emphasis on Christ's prior gift is the controlling theological order: divine generosity grounds the mission, and the workers' action is derivative, accountable, and directed toward persons Jesus has already seen.

The Twelve are individually named before they are sent. Their mission is corporate and personal, authorized yet bounded. The initial restriction to Israel belongs to this stage of Matthew's narrative and to the covenant history within which the mission begins. The final command to give without cost recalls that authority, healing, and proclamation are received gifts. It rules out ownership and exploitation, not the wider New Testament's lawful support of ministers.

Chrysostom makes the anti-mercenary point without denying support. The apostles must not commercialize healings or teaching because the capacity was itself received as gift; later instructions about food and lodging still recognize that workers may be sustained. The distinction matters wherever “without cost” is used either to excuse spiritual profiteering or to deny ministers ordinary justice. The mission's initial destination likewise belongs to salvation history: Jesus attends first to the people among whom covenant promises have been given, and Matthew's Gospel later narrates the universal commission.

\subsection*{The independent apostolic foundation\enspace\elementref{Romans}}

Romans 5 does not function as a hidden allegory of Exodus or the Twelve. Paul intensifies humanity's condition: weak, ungodly, sinful, and hostile. Christ's death does not answer human worthiness but creates reconciliation. Chrysostom's direct exposition emphasizes the sequence from enemies reconciled through death to a stronger hope of salvation through Christ's living power.

Aquinas's \emph{Commentary on Romans} V, lectio 2, reads the appointed movement as an argument from the harder gift to the assured completion: if God reconciled enemies through the Son's death, salvation through his life is not a weaker hope. He distinguishes the successive descriptions of the beneficiaries without turning them into classes of people whom the Church may rank. Augustine, arguing against Pelagian accounts in \emph{Against Two Letters of the Pelagians} IV.8, presses Romans 5:10--12 to show that reconciliation presupposes an enmity caused by sin and a grace not generated by unaided nature. His polemical setting is later and narrower than Paul's letter, but the use is direct, and it bars any reading of the passage as a celebration of latent human sufficiency.

The final phrase about rejoicing in God through Christ prevents reconciliation from becoming only a legal change described from outside. Gift establishes a relationship in which the reconciled may glory in the giver. Yet the appointed verses stop before Paul's Adam--Christ comparison, and their death-to-life argument and gratuitous reconciliation stand without Romans 5:12--21.

Chrysostom and Aquinas also preserve the argument's temporal force. Paul does not say merely that God might love enemies after they improve: the death occurs while the beneficiaries are weak and hostile. The later “much more” is not confidence in their unaided constancy but confidence in the living Christ whose reconciling act has already crossed the greater distance. Augustine's anti-Pelagian use makes the same priority polemically, though Paul's sequence does not carry every contour of that controversy. Missioners formed by this reading cannot divide the field into those already worthy of grace and those beyond it. They themselves stand inside the category Christ's act has changed.

This strand changes the spiritual posture of mission. Workers are first beneficiaries. “Without cost” can therefore be heard under Romans' priority of gift as a source-grounded synthesis: the Church gives what she did not originate, to persons she may not classify as deserving or undeserving.

\subsection*{Kingdom acclamation\enspace\elementref{Acclamation}}

Mark's verse is acclamatory: it announces fulfilled time, nearness of the Kingdom, conversion, and faith in order to welcome Christ about to speak. Its role ends with that welcome; it is not another homily before the Gospel. Its Kingdom language nonetheless clarifies the mission's object: the Twelve announce divine reign, not themselves.

Tatian's second-century \emph{Diatessaron}, section 5, receives the opening proclamation in a Gospel harmony: the beginning of Jesus' ministry is marked by fulfilled time, the nearness of God's reign, repentance, and belief. That is early direct reception of the saying, although the harmony does not preserve Mark's literary sequence as a modern commentary would. Its compactness suits the liturgical role. The acclamation does not supply a second account of the Twelve; it places their mission under the prior proclamation of Jesus, whose Kingdom calls hearers to conversion before it sends them to serve.

\subsection*{Gift and sacramental unity\enspace\elementref{Prayer over Offerings, Communion, Prayer after Communion}}

The Prayer over the Offerings compares the created nourishment of bodily nature with sacramental renewal, then asks aid for both body and mind. This refuses a spirituality that treats embodiment as disposable. It also prevents an automatic promise of present bodily cure: the prayer concerns God's sustaining and sanctifying gift within the Church's sacramental horizon.

Communion option A continues Psalm 27's movement from plea to dwelling. Option B takes Christ's prayer that his disciples be one. The Prayer after Communion makes unity more than sentiment: reception signifies it and petitions its effect. Division cannot be healed by suppressing truth or justice, and sacramental reception is not magic; nevertheless, unity is God's gift before it becomes the Church's ethical work.

For option A, Augustine's psalm exposition makes the sought dwelling a stable orientation of desire, not an escape into sacred space that leaves conduct untouched. For option B, Chrysostom's \emph{Homily 81 on John} reads John 17:11 within Jesus' departure: the disciples remain in the world, Christ is glorified in them, and their preservation depends on the Father. Augustine's \emph{Tractate 110 on John} directly treats the prayer's later unity language and carefully refuses to make created ecclesial unity identical with the consubstantial unity of Father and Son. The Church's unity is real gift and participation, not absorption or uniformity. These receptions belong only to their respective appointed alternatives; which antiphon the dated occurrence used is unresolved.

The appointed John verse asks preservation in the Father's name before it states the purpose “that they may be one.” Unity is therefore received under divine keeping, not manufactured by minimizing difference or conflict. Chrysostom's attention to disciples remaining in the world also excludes a merely inward unity detached from mission and danger. Augustine's later-verse distinction supplies a doctrinal boundary: the analogy to Father and Son makes ecclesial communion demanding, but creatures do not become one substance. The Prayer after Communion can ask for unity as sacramental effect precisely because the Gospel branch has first located it in God's action.

\Needspace{16\baselineskip}
\section*{The Propers: Notable and Quotable}

{\small
\begin{itemize}
\item \textbf{“I bore you up on eagle wings” --- an airlift becomes the metaphor.\enspace\elementref{First Reading}} Israel's 1949--1950 evacuation of Yemenite Jews was also called \emph{Operation On Wings of Eagles}. The aircraft made Exodus' remembered rescue startlingly concrete: a scriptural image supplied the public name for a modern journey to Israel.

\item \textbf{“The names of the twelve apostles are these” --- except at an Australian landmark.\enspace\elementref{Gospel}} Victoria's limestone stacks were advertised as the “Twelve Apostles” by the 1920s. A commissioned heritage study later contradicts that evidence by placing the probable renaming in the 1950s, so the exact date remains unresolved. Parks Victoria now dryly notes that there were never actually twelve. Matthew's apostolic count became a memorable tourist brand whose joke is its geographical inaccuracy.

\item \textbf{“Without cost you have received; without cost you are to give” --- a software-license epigraph.\enspace\elementref{Gospel}} The Bible SuperSearch client places the King James form, “freely ye have received, freely give,” beneath its GPL notice and immediately before detailed commercial-use restrictions. Missionary instruction becomes a slogan for distributable code; the adjacent price and licensing qualifications make “free” do unexpectedly technical---and faintly ironic---work.
\end{itemize}
}

\section*{The Propers: Interpretive Possibilities}

\begin{studybox}{Editorial status}
The following exploratory proposals were developed in an AI-assisted editorial process. They are attributed to none of the cited authorities and claim no historical compositional intent. Each states its anchors and limits.
\end{studybox}

\Needspace{10\baselineskip}
\noindent\textbf{Election takes the form of a gift to former enemies.\enspace\elementref{First Reading; Second Reading; Gospel}}
Exodus calls a people carried, treasured, and priestly; Romans locates Christ's gift among the weak, sinful, and hostile; Matthew sends the Twelve to give freely what they received. In sequence, election circulates mercy: those gathered by grace mediate an unowned gift to people unable to qualify themselves first. Separate reading can isolate covenant identity, enemy-reconciliation, and mission; together they test privilege by received mercy. The fruit is an anti-triumphalist vocation ordered to God's universal purpose. An official homily on the exact Exodus and Gospel readings supplies precedent for that ordering, though Romans is semi-continuous. The proposal retains Israel's election, covenant distinctions, the initial mission to Israel, and prudent boundaries toward active harm.

\Needspace{10\baselineskip}
\noindent\textbf{The flock becomes laborers without ceasing to be flock.\enspace\elementref{Responsorial Psalm; Gospel}}
Psalm 100 has the assembly confess that God made it and it remains his flock; Matthew sees a shepherdless crowd, commands prayer for workers, then names and authorizes the Twelve. The juxtaposition changes grammatical role without transferring ownership: receivers of shepherding become harvest agents while God remains maker, shepherd, and field's Lord. Separate reading misses dependence converted into responsible agency and can divide passive laity from self-originating experts. The fruit is a baptismal culture in which care equips care and every minister remains Christ's sheep. No checked source joined Psalm 100 to this exact transition; Chrysostom identifies the prayed-for laborers with the Twelve. Flock and harvest are not interchangeable, and the proposal neither equates every office nor licenses unformed or unauthorized ministry.

\Needspace{10\baselineskip}
\noindent\textbf{The prayer for laborers begins to answer itself.\enspace\elementref{Collect; Gospel}}
The Collect asks grace to govern willing and action; in Matthew the disciples ask for workers, then are named, empowered, and sent. Petition thus forms a vocational feedback loop: rather than recruit an absent third party, it places petitioners where grace can turn compassion into consent and action. Reading prayer and commission separately misses a community becoming part of its answer. Chrysostom observes that Christ commands prayer and then ordains the Twelve; modern official teaching calls prayer the first service to vocation. Extending that precedent through the Collect yields planning that begins dependently yet expects costly availability. It neither gives every petitioner the same call, replaces discernment and formation, nor confines “laborers” to ordained ministry.

\Needspace{10\baselineskip}
\noindent\textbf{Mission is enclosed by one undivided desire.\enspace\elementref{Entrance; Gospel; Communion A}}
The Entrance takes Psalm 27's cry for help; Communion A, if appointed, takes its desire to dwell in God's house; between them Matthew sends the Twelve into towns, bodies, and households. This same-psalm frame may make mission an outward-and-homeward movement sustained by one desire: communion with God sought for sender and recipient. Separate reading notices repetition and commission but misses that rhythm. Its fruit is a test against activism, institutional self-enlargement, and withdrawal disguised as contemplation. No exact precedent was located. Communion A is unresolved, dwelling is not withdrawal, and the frame proves neither that the antiphon was chosen nor that the Missal historically designed the arc.

\Needspace{10\baselineskip}
\noindent\textbf{Unity can hold a mission that multiplies names and places.\enspace\elementref{Gospel; Communion B; Prayer after Communion}}
Matthew names and sends distinct apostles; Communion B, if appointed, asks the Father to keep disciples in the divine name as one; the final prayer asks sacramental participation to bring unity. Divine keeping, not standardization, lets differentiated persons disperse without becoming rival proprietors because communion comes from their sender. Separate reading can reduce the Twelve to a list and unity to internal harmony, missing communion as the condition of noncompetitive plurality. Official teaching on apostolic collegiality and Augustine on John 17 provide precedent, though not for the whole chain. The fruit is accountable decentralization and truthful collaboration. Communion B is unresolved; unity is not divine consubstantiality and cannot suppress evidence, conflict, charisms, legitimate authority, or harmed persons' voices.
